\documentclass[11pt]{article}

\usepackage{amsmath,amssymb}
\usepackage{xurl}
\usepackage[hidelinks]{hyperref}
\setlength{\emergencystretch}{2em}

% Shared commands for notes in docs/paper-gaps/.
% Notation follows the LDT paper (references/ldt-paper/).

\newcommand{\Lean}{\textsc{Lean}}
\newcommand{\leanid}[1]{\textsf{#1}}
\newcommand{\ghissue}[1]{\href{https://github.com/LionSR/MIPStarRE/issues/#1}{GitHub issue \##1}}
\newcommand{\ghpr}[1]{\href{https://github.com/LionSR/MIPStarRE/pull/#1}{GitHub PR \##1}}

% --- Paper-compatible notation ---
\newcommand{\F}{\mathbb{F}}
\newcommand{\N}{\mathbb{N}}
\newcommand{\C}{\mathbb{C}}
\newcommand{\tr}{\operatorname{tr}}
\newcommand{\ot}{\otimes}
\newcommand{\E}{\mathop{\bf E\/}}
\newcommand{\bx}{\mathbf{x}}
\newcommand{\by}{\mathbf{y}}
\newcommand{\bu}{\mathbf{u}}
\newcommand{\bv}{\mathbf{v}}

% Approximation relation (paper uses \approx_\eps and \simeq_\zeta)
\newcommand{\approxeps}{\approx_{\varepsilon}}
\newcommand{\simgeq}{\simeq_{\zeta}}

% Polynomial spaces (paper uses \polyfunc, \polysub, \polymeas)
\newcommand{\polyfunc}[3]{\operatorname{Poly}(#1,#2,#3)}
\newcommand{\polysub}[3]{\operatorname{PolySub}(#1,#2,#3)}
\newcommand{\polymeas}[3]{\operatorname{PolyMeas}(#1,#2,#3)}


\title{Documenting Possible Gaps, Counterexamples, and Formalization Deviations}
\author{}
\date{2026-04-29}

\begin{document}
\maketitle

\section{Purpose}

When a formalization relies on a cited mathematical source, the cited source
should remain intact.  It is a source, not an errata file.  If a possible gap,
missing hypothesis, incorrect numerical constant, or other mathematical
discrepancy is found, the matter should be recorded in a separate LaTeX note.
In this repository such notes live under \path{docs/paper-gaps/}.  The same
principle applies when the formal development has already diverged from the
cited source, for instance by changing a statement so that it compiles, or by
using a different mathematical route because an available formal-library theorem
has a slightly different form.  Such a note is a mathematical note first, and a
record of the formalization second: it should be readable by a mathematician or
mathematical physicist who has not followed the issue discussion.

The note should isolate the mathematical mechanism.  It should say what the
cited source asserts, what the cited results actually imply, where a term or
hypothesis is lost, and what statement the formal development proves instead.
If the formal proof takes a different route, the note should explain the
replacement argument and the hypotheses on which it rests.  This is especially
important when the proof uses a theorem, construction, or convention from
Mathlib which is not part of the cited argument.  In that case the note should
state the Mathlib result in ordinary mathematical language, explain how it is
specialized to the notation of the paper, and compare its hypotheses and
conclusion with the step it replaces.  It should not merely say that \Lean{} did
not accept a proof, that the modified proof compiles, or that a named library
lemma was applied.

The documents in this directory should themselves follow this standard.  A
policy or template which reads like a checklist will tend to produce checklist
notes; the examples and templates should therefore be written as mathematical
prose.

\section{When the Investigation Starts}

The investigation starts when the obstruction appears to be mathematical rather
than merely formal.  No final judgement is needed at this stage.  It is enough
that, after reading the relevant source and attempting the corresponding formal
statement, the proof seems to require a different scalar, an additional
hypothesis, a different tensor placement, a distinction between two objects which
the prose identifies, or an auxiliary assertion not supplied by the cited
argument.

The same investigation should also start when such a deviation is discovered
after the fact.  If the existing formal development or informal blueprint has
already changed the statement, strengthened a hypothesis, weakened a conclusion,
inserted an unadvertised lemma, or replaced the cited proof by a different
formal-library argument, the mathematical change should be documented even if all
relevant files already compile.

At this early stage the purpose is to preserve the question while it is still
local.  The note may later conclude that the matter is only notational, that an
intermediate estimate in the cited source needs a harmless correction, that a
missing hypothesis is genuinely needed, or that no discrepancy remains.  Thus
the trigger is not the severity of the eventual verdict, but the appearance of a
real mismatch between the cited reasoning and the statement that can be
formalized.

A failed tactic search, a missing formal lemma, or an inconvenient coercion is
not enough.  The obstruction must survive translation back into ordinary
mathematics.

\section{When to Write a Note}

A note should be written once the source argument and the formal proof no longer
match as mathematical arguments.  This includes, for instance, a numerical
estimate whose constant does not follow from the cited lemma, a displayed
conclusion which uses an unstated hypothesis, or a step in which two objects are
identified although they must remain distinct.  In analytic, algebraic,
probabilistic, or quantum arguments, such distinctions may concern the domain of
a map, the placement of a tensor factor, the type of an outcome, the distinction
between a structure and a substructure, or the scalar carried by an approximation
relation.

A note is also appropriate when the proof becomes formal only after changing an
intermediate scalar, weakening a conclusion, strengthening a hypothesis, or
inserting an additional lemma which is not simply an expansion of the cited
argument.  These are mathematical changes, even when they are harmless for the
final theorem.  It is likewise appropriate when an available formal-library
theorem proves the needed result under a somewhat different formulation from the
cited source.  In that case the note should be deliberately pedagogical.  It
should explain, in ordinary mathematical language, what the replacement result
says, how it is specialized, which hypotheses are new or differently packaged,
and why the resulting statement is a legitimate substitute for the step in the
cited argument.  If the issue remains open after reasonable attempts, the note
should say what has been checked and what remains unknown.

\section{Counterexamples}

If the cited assertion appears to be false, the note should not stop at saying
that the proof cannot be formalized.  It should give the counterexample, or the
smallest relevant part of it, in ordinary mathematical notation.  The example
should be accompanied by enough verification that the reader can see both sides:
the hypotheses of the cited assertion hold, and its conclusion fails.

When feasible, the counterexample should also be formalized in \Lean{}.  This is
especially valuable when the example is finite, algebraic, computational, or
otherwise small enough that the failure can be checked without reproducing the
whole surrounding theory.  The note should then state the mathematical content
of the checked counterexample and identify the corresponding Lean declaration in
a footnote.  The Lean file is not a substitute for the prose: the prose should
still explain why the example refutes the cited assertion.

\section{Revision of a Note}

A paper-gap note may be provisional.  It may begin as soon as the mathematical
obstruction has been isolated well enough to state the question, and it may be
revised as the calculation, the blueprint, or the corresponding \Lean{}
declarations become clearer.  The first version need not know the final
classification, but it should say plainly which mathematical point is known and
which point remains unsettled.

Later revisions should not make the note into a chronological log.  They should
rewrite the document so that it remains a coherent account of the present
mathematical understanding.  If the verdict changes, the conclusion should state
the current verdict; the linked issue or pull request can preserve the earlier
discussion.

\section{Style and Notation}

\subsection{Notation must match the cited source}

Every paper-gap note must use the notation of the cited source and must not
silently replace it with a different convention.  If the source writes
$\F_q$ for the finite field, the note writes $\F_q$, not $\mathbb{F}_q$.
If the source uses $\tr$ for the trace and $\ot$ for the tensor product,
the note does the same.

\begin{itemize}
  \item Good (source uses \verb|\F_q|): ``We work over the finite field $\F_q$.''
  \item Bad (silently replaced): ``We work over the finite field $\mathbb{F}_q$.''
  \item Good (source uses \verb|\ot|): ``The operator $A \ot I$ acts on the
        first register.''
  \item Bad: ``The operator $A \otimes I$ acts on the first register.''
\end{itemize}

The shared command file \path{docs/paper-gaps/command.tex} should define
the macros used by the source with the same expansions, so that a note
compiled standalone produces notation identical to the source.  A note must
not redefine any command already defined in \path{command.tex} or introduce
a competing convention for the same symbol.

When a paper-gap directory serves a particular source (as the present one
serves \cite{Ji2020LowIndividualDegree}), the \path{command.tex} for that
directory is the single point of truth for notation.  If a macro is missing,
add it there; do not add local definitions in individual notes.  If the
source changes its notation in a later version, update \path{command.tex}
and rebuild the notes; do not fork the notation across files.

\subsection{Discipline on \path{command.tex}}

The file \path{command.tex} should contain only commands that appear in the
cited source or are genuinely needed by every note (such as \verb|\Lean| and
\verb|\leanid|).  Do not add a command to \path{command.tex} unless the
source paper defines it, and do not change the expansion of an existing
command to differ from the source.  If the source paper writes
\verb|\newcommand{\bx}{\boldsymbol{x}}|, then \path{command.tex} must do the
same; do not replace it with \verb|\mathbf{x}| or \verb|\mathbb{x}|.  A
mathematician who knows the source paper's notation should see the same
symbols in a paper-gap note without having to check a separate style file.

\subsection{Voice}

A paper-gap note should write in the same voice as the source it examines.
For most mathematical papers this means the first-person plural and direct
declarative sentences.  Do not use decorative language, hedging phrases,
or AI-generated jargon.  Name the mathematical object or operation in
ordinary mathematical language.

A mathematician reading the note should feel that it could have been written
by a colleague, not by a project management tool.  The note is part of the
mathematical record, not part of the project tracker.

\begin{itemize}
  \item Good: ``We examine the step at line~167.  The proof applies the
        triangle inequality to the self-consistency estimate and the
        closeness bound, obtaining $\zeta_1$ exactly.''
  \item Bad: ``This note concerns the line~167 projectivization transport
        step.  A faithfulness audit identified a scalar/tensor bridge
        divergence requiring the P-level match-mass monotonicity
        invariant.''
  \item Good: ``The constant in the paper is $6$, but the cited lemma gives
        $\sqrt{32}$.  Since $\sqrt{32} \le 6$, the paper's bound still
        holds after rounding up.''
  \item Bad: ``The Lean formalization uses $\sqrt{32}$ while the paper
        states $6$.  This discrepancy arises because the formal proof
        cannot synthesize the rounding step automatically.''
\end{itemize}

Notice that the bad examples read like issue tracker prose: they name
formal declarations, use project-internal jargon (``faithfulness audit'',
``P-level''), and describe the formalization process rather than the
mathematics.  The good examples state the mathematical facts directly and
let the reader judge.

\subsection{Lean identifiers belong in footnotes}

Lean declaration names, type names, file paths, and GitHub issue numbers
belong in footnotes, not in the body text.  The body must read as pure
mathematical prose.

\begin{itemize}
  \item Good: ``The formal development represents a distribution as a finite
        weighted sum.\footnote{The type is called \leanid{Distribution} in
        \path{MIPStarRE/LDT/Basic/Distribution.lean}.}''
  \item Bad: ``The code defines a \leanid{Distribution} type carrying a
        \leanid{Finset} support with an \leanid{IsProbability} predicate.''
  \item Good: ``The lemma is proved in the formalization.\footnote{The
        declaration is \leanid{commDataProcessedG} in
        \path{ProcessedG.lean}.}''
  \item Bad: ``The Lean lemma \leanid{commDataProcessedG} replicates the
        same ten-step chain.''
\end{itemize}

When a Lean identifier must appear in a footnote, use the semantic command
\verb|\leanid{Name}| (which expands to \leanid{Name}), not the generic
\verb|\texttt{Name}|.  GitHub issues and pull requests should be cited with
the shared commands \verb|\ghissue{N}| and \verb|\ghpr{N}|.

\subsection{Titles}

The title should name the mathematical topic, not the formal declaration under
discussion.  For example, ``Why the transport carries an extra
$2\sqrt{\zeta}$'' rather than ``The Scalar/Tensor Bridge in
\leanid{fullSliceABABAvg}''.  More generally:

\begin{itemize}
  \item Good: ``A missing hypothesis in the contraction estimate''
  \item Bad: ``Issue 427: \leanid{errorCascadeBound} needs stronger input''
  \item Good: ``The constant in the variance bound''
  \item Bad: ``Fixing the $10000$ vs $20000$ mismatch in
        \leanid{globalVarianceBound}''
\end{itemize}

The title tells the reader what mathematical question the note addresses.
It should not read like a GitHub issue title.

\subsection{The verdict}

The concluding section of a note is called ``Conclusion''.  It should open
directly with the verdict, stated as a plain mathematical assertion, and then
justify it in a short paragraph.  Do not preface the verdict with a label.

\begin{itemize}
  \item Good: ``There is no mathematical discrepancy.  The formalization
        separates the prime-power hypothesis on $q$ from the use of a finite
        field of order $q$, and the two are linked by a canonical
        isomorphism.  Every hypothesis of the paper is preserved.''
  \item Bad: ``\textbf{Verdict:} no discrepancy.''
\end{itemize}

The section heading already tells the reader that this is the conclusion.
The first sentence should state the verdict; the rest of the paragraph
should give the reasoning.  A verdict that is a single sentence without
justification reads as a checklist item, not as a mathematical conclusion.

\section{Form of the Note}

Each note should stand on its own.  It should give the date, cite the relevant
source article, and provide enough traceability for the reader to locate the
original passage and the project discussion.  GitHub issues and pull requests may
be linked with \verb|\href|, and source files, line ranges, and labels should be
recorded when they help the reader check the claim.  These references should
serve the exposition rather than govern it: the main text should be organized
around the mathematical assertion and its proof.

The note should introduce its notation before using it.  Symbols for
measurements, tensor factors, approximation relations, and error parameters
should be named in prose.  If an approximation relation is state-dependent, the
state should be named, or its defining expression should be given when this can
be done briefly.  Only after the notation is in place should the note state the
assertion in the cited source, using a bibliographic citation such as
\verb|\cite[Section~...]{...}| when the section or chapter is known.

Only the relevant hypotheses should be reproduced.  If the question is a scalar
substitution, a missing term, or a missing hypothesis, the note should not quote
a long unrelated part of the proof.  It should give the cited lemma or
proposition in the special case actually used, then display the calculation or
logical implication in which the discrepancy appears.

Displayed equations should be used with restraint.  They are appropriate for the
source estimate being checked, the cited estimate in the form actually used, and
the calculation where a term or hypothesis appears or disappears.  When a
displayed equation is identified by a source label, the label should appear on
the right of the display, while the surrounding prose refers to the mathematical
assertion itself.  The standard \texttt{amsmath} command \verb|\tag| is the
right mechanism for this; a separate local convention should be introduced only
if the standard mechanism is genuinely insufficient.  For a multi-line aligned
display, tag the individual lines that carry the source information rather than
forcing a long combined label below the equations.  Short
scalar definitions and elementary inequalities should remain inline when the
sentence reads better that way.  As in ordinary mathematical writing, a displayed
formula is part of the surrounding sentence and should be punctuated accordingly.

Bibliographic entries used only by these notes should live in
\path{docs/paper-gaps/references.bib}.  A new note should ordinarily begin from
\path{docs/paper-gaps/template.tex}, which records the common LaTeX preamble,
the local bibliography, and a model for right-hand source labels in displayed
equations.

The note should then compare the cited source with the informal blueprint, if
there is one, and with the formal statement.  This comparison should be written
in ordinary mathematical prose.  Formal declaration names and source paths should
be supplied only to identify the exact formal statements, and normally belong in
footnotes.

The conclusion should give a verdict.  It may say that there is no discrepancy,
that the matter is only notational, that an intermediate estimate needs a
harmless correction, that a hypothesis is missing, that there is a theorem-level
gap, or that the question remains open.  The classification should reflect the
mathematics, not the amount of formal work required.

\section{Decision Principle}

The formal development should not silently replace the cited assertion by a
different one.  If the formal statement differs from the cited source, the
reason for the difference should be written down.  Conversely, the note should
not overstate the conclusion.  A wrong intermediate scalar is not a fatal gap in
the theorem if the corrected scalar is bounded by the error parameter appearing
in the theorem statement.

\end{document}
